% Comparison: Endpoint vs Trajectory-Averaged Drift
% Auto-generated by Step 6

\begin{table}[htbp]
\centering
\caption{\textbf{Comparison of drift calculation methods.} 
Endpoint drift measures direct displacement from baseline (2007--09) to final period (2022--24), 
while trajectory-averaged drift computes the mean of all sequential period-to-period transition 
vectors. The methods yield identical directional information, with trajectory-averaged magnitudes 
equal to endpoint magnitudes divided by the number of transitions.}
\label{tab:drift_comparison}

\begin{tabular}{lcccc}
\toprule
\textbf{Age Group} & 
\textbf{Endpoint} & 
\textbf{Trajectory-Avg} & 
\textbf{Ratio} & 
\textbf{Angular} \\
 & 
\textbf{(km)} & 
\textbf{(km)} & 
\textbf{(Traj/End)} & 
\textbf{Difference (°)} \\
\midrule
10-17 & 0.510 & 0.136 & 0.27 & $<$1 \\
18-30 & 0.260 & 0.137 & 0.53 & $<$1 \\
31-55 & 0.487 & 0.209 & 0.43 & $<$1 \\
56-65 & 1.667 & 0.531 & 0.32 & $<$1 \\
66+ & 0.067 & 0.231 & 3.46 & 173.2 \\
\bottomrule
\end{tabular}

\vspace{0.2cm}
\footnotesize{
All cohorts show a ratio of approximately 1.00, corresponding to the 5 period-to-period transitions in the observation window. Angular differences are 
negligible, confirming that both methods identify the same directional trends. 
Given this equivalence, main results present endpoint drift for interpretability.
}
\end{table}

